\paragraph{Unfolding loops}
In languages like \texttt{C} (and by extension \texttt{C++}) and \texttt{Java}, unfolding a loop leads to non-equivalent code:

\begin{multicols}{2}
    \begin{code}{c}
        int i = 0;
        while ( i < 2 ) {


            while ( i < 1 )
                if ( i == 0 ) break;

            i++;
        }
        printf( "i = %d", i );
    \end{code}
    Prints \shade{gray}{\texttt{i = 2}}
    \begin{code}{c}
        int i = 0;
        while ( i < 2 ) {
            if ( i == 0 ) {
                if ( i == 0 ) break;
                while ( i < 1 )
                    if ( i == 0 ) break;
            }
            i++;
        }
        printf( "i = %d", i );
    \end{code}
    Prints \shade{gray}{\texttt{i = 0}}
\end{multicols}

In IMP however, this kind of action leads to equivalence:
\[
    \forall b, s.(\texttt{while b do s end} \; \simeq \; \texttt{if b then s; while b do s end end})
\]
So what we have to prove is this:
\[
    \forall b, s, \sigma, \sigma'.(\vdash \langle \texttt{while b do s end}, \sigma \rangle \; \simeq \;
    \vdash \langle \texttt{if b then s; while b do s end end}, \sigma \rangle \rightarrow \sigma')
\]

A proof idea is to show equivalence by proving the implication in both directions. For each of them, we show that there is a derivation tree.

\inlineproof available in the lecture slides for Formal Methods, Slides 78 - 80 (Slide Deck 3, pages 23 - 25)
